\section*{Introduction}

This report presents the design, implementation, simulation, and evaluation of an
agentic system for zero-touch connectivity and traffic engineering. The main design
question was not whether a large language model (LLM) could generate routes from
scratch, but rather which role an LLM should play inside a routing system that must
remain safe, responsive, and compliant with network constraints.

Traditional graph algorithms are faster and more reliable than an LLM for tasks such
as verifying adjacency, computing paths, and enforcing hard capacity constraints.
Conversely, an LLM can be useful as a supervisory component that interprets a compact
network situation and selects a high-level routing objective. The final architecture
therefore separates strategic reasoning from deterministic execution: the LLM chooses
the objective, while deterministic tools generate, allocate, and validate the routes.

The implementation uses LangGraph to make this decision process explicit. It also
uses asynchronous, event-driven LLM calls so that model inference does not block the
simulation. The final system is evaluated against three heuristic baselines and with
several remote and local language models.


\section{Design Process}

The final architecture was reached through several prototypes. Although more than
three code versions were implemented, they can be grouped into three main design
families. Each family explored a different answer to the central question: how can an
LLM influence routing without replacing the deterministic mechanisms that are already
well suited to path computation?

\subsection{First design: per-flow path selection}

The first design used the LLM as a selector of individual paths. A perception node
collected the network state, a deterministic function generated candidate routes for
each flow, and the LLM selected one candidate. The selected route was then validated
before installation.

This approach reduced hallucination risk because the LLM could only choose among
existing graph paths. It also gave the model direct influence over every routing
decision. However, its limitations became evident as the simulation progressed:

\begin{itemize}
    \item The number of active flows increased continuously, which enlarged the prompt
    and increased token consumption.
    \item Per-flow reasoning required too many model decisions and produced excessive
    inference latency.
    \item By the time a response arrived, the network state could already have changed,
    making the decision stale.
    \item Processing flows independently did not provide a reliable global view of
    shared capacity. Two individually feasible routes could compete for the same link
    and become infeasible when installed together.
\end{itemize}

The main lesson from this version was that an LLM should not operate as a slow
replacement for a shortest-path algorithm, especially when the number of flows grows
over time.

\subsection{Second design: global routing policies}

The second design changed the role of the LLM from per-flow path selection to global
policy generation. Instead of selecting every route, the model observed an aggregated
description of the network and selected global weights or priorities that were used by
a deterministic routing algorithm. One LLM decision could therefore influence many
flows.

This version reduced prompt size and model-call frequency. It also ensured that valid
graph paths were generated by deterministic code. Nevertheless, a set of global
weights has an important limitation: it can express a general preference for latency,
utilization, or capacity, but it cannot directly target a specific bottleneck. For
example, increasing the global penalty for utilization affects every link, even if the
real problem is concentrated on only one link.

Two additional ideas were evaluated during this stage. First, a fixed capacity reserve
was introduced to preserve headroom. In practice, an aggressive reserve made the
allocator unnecessarily conservative and increased the number of rejected flows. This
motivated the use of a shared, prospective capacity ledger instead of subtracting the
same reserve from every link. Second, short-term memory was added so that the LLM could
observe previous choices and outcomes. It did not consistently improve performance:
the topology and offered traffic changed quickly, and the system lacked a sufficiently
clean causal link between one policy decision and a later outcome. The memory therefore
added context and latency while sometimes preserving stale conclusions.

The global-policy approach was safer and more scalable than per-flow selection, but it
still required a more precise mechanism for deciding when the model was necessary and
for applying its decision consistently across all flows.

\subsection{Third design: hybrid fast path and LLM path}

The third design introduced triage. The agent first classified the network situation.
Routine cases were sent to a deterministic fast path, while conflicts and failures
were sent to the LLM. This substantially reduced unnecessary calls: a simple direct
route did not require language-model reasoning.

This architecture improved responsiveness, but early versions processed fast-path and
LLM-managed flows separately. That separation created a consistency problem. The LLM
could select a strategy under the assumption that capacity was available, while the
fast path had already assigned that capacity to other flows. In addition, most flows
were resolved deterministically before the LLM response arrived, so the model had only
a weak practical role.

Several intermediate prototypes explored caching, Top-$K$ affected flows, component
allocation, pressure thresholds, and different candidate-ranking functions. Reusing a
cached strategy solely because it had been selected before proved too coarse: two
events with the same label could involve different flows and links. The final design
therefore reuses an active plan only while a compact event signature remains unchanged,
and it replaces separate allocation with one joint allocator that sees all active
flows.

\subsection{Key lessons}

The final design was guided by three main lessons:

\begin{itemize}
    \item \textbf{LLM role:} the model should select a global routing objective, not
    invent topology paths or independently route every flow. Deterministic tools remain
    responsible for path generation and hard constraints.

    \item \textbf{Event-driven inference:} the LLM should not be called at every
    simulation cycle. It is called when the network event changes meaningfully, while
    a valid active plan can be reused between related cycles.

    \item \textbf{Global allocation:} all active flows must share the same capacity
    view. A joint allocator applies both routine and strategy-driven decisions using a
    single reservation ledger, preventing the fast path and LLM path from allocating
    the same capacity independently.
\end{itemize}


\section{Final Architecture}

The final implementation is a LangGraph workflow with the following nodes:

\begin{center}
\texttt{perceive} $\rightarrow$ \texttt{classify} $\rightarrow$
$\{\texttt{fast\_path},\ \texttt{plan\_llm}\}$ $\rightarrow$
\texttt{route\_with\_strategy} $\rightarrow$ \texttt{validate}
$\rightarrow$ \texttt{record} $\rightarrow$ \texttt{END}.
\end{center}

The conditional edge after classification is part of the agent policy: deterministic
routing is not an accidental fallback, but an explicitly selected tool for routine
cases. Complex cases require a valid LLM plan before strategy-based allocation is
performed.

\subsection{Perceive}

The \texttt{perceive} node collects the topology and flow-table snapshots once per
decision cycle. The topology snapshot contains nodes, links, capacities, latencies,
loss values, observed utilization, and link status. The flow-table snapshot contains
previously installed paths, allowing the agent to preserve routes that remain valid
instead of recalculating the complete network unnecessarily.

This node also checks whether an asynchronous LLM request has completed. It never waits
actively for the response. A completed and valid response becomes the active plan; an
invalid or empty response is rejected and does not silently install a default strategy.
This distinction is important because the absence of an LLM answer must not bypass the
agentic decision requirement.

\subsection{Classify}

The \texttt{classify} node identifies whether the current situation is
\texttt{trivial}, a \texttt{conflict}, or a \texttt{failure}. For every flow, it
checks whether the direct source--destination link exists, is operational, has enough
nominal capacity, and satisfies the flow's maximum-latency SLA. A direct route being
insufficient does not mean that the flow is impossible; it means that a multi-hop
candidate or a coordinated capacity decision may be necessary.

The node also computes prospective link pressure. For a link $e$, pressure is defined
as

\begin{equation}
    P_e = \frac{\sum_{f \in F_e} d_f}{C_e},
\end{equation}

where $d_f$ is the demand of a flow expected to use link $e$, $F_e$ is the set of such
flows, and $C_e$ is the link capacity. The maximum link pressure summarizes whether
individually valid routes are likely to conflict when considered together. Pressure is
prospective, whereas observed utilization describes traffic already present in the
network.

The resulting classification is:

\begin{itemize}
    \item \texttt{trivial}: direct routing is sufficient and prospective pressure is
    below the configured threshold;
    \item \texttt{conflict}: at least one flow needs a multi-hop route, or prospective
    pressure indicates shared-capacity contention;
    \item \texttt{failure}: at least one network link is down.
\end{itemize}

After classification, the node creates an \texttt{event\_key}: a compact signature
containing the situation type, bucketed numbers of active and affected flows, affected
demand, failed-link identities, utilization and pressure bands, and the highest-pressure
links. The LLM is requested only when this signature changes meaningfully. This is plan
persistence rather than long-term memory: the active strategy is reused while the same
network event remains in effect.

\subsection{Fast Path}

The \texttt{fast\_path} node is selected only for trivial situations. It assigns
feasible direct routes through a shared bidirectional capacity ledger. If a 10 Mbps link
has already reserved 6 Mbps during the current allocation cycle, later flows see only
4 Mbps of remaining capacity. Therefore, flows are not accepted independently when
their aggregate demand would overload a common link.

The fast path is intentionally narrow. It handles only cases that require no strategic
choice; it does not hide an LLM failure or assign complex multi-hop routes using an
implicit default policy.

\subsection{Plan with LLM}

The \texttt{plan\_llm} node is used for conflicts and failures. The model receives an
aggregated network description rather than raw topology and flow lists. The prompt
contains the event type, active and affected-flow counts, failed and high-pressure links,
utilization and pressure levels, and compact candidate information for the Top-$K$
affected flows. The Top-$K$ flows are selected by demand so that the prompt concentrates
on the requests with the largest immediate capacity impact. The allocator later applies
the selected objective to all active flows, not only to those shown to the model.

The LLM must select exactly one objective from a closed set:

\begin{itemize}
    \item \texttt{maximize\_acceptance}: prioritize flows with scarce alternatives and
    smaller demands, with the objective of admitting more requests;
    \item \texttt{minimize\_delay}: prioritize tight latency SLAs and rank candidates
    by their end-to-end latency;
    \item \texttt{protect\_capacity}: avoid candidates with high projected utilization
    and preserve capacity headroom;
    \item \texttt{stable\_recovery}: preserve valid installed routes whenever possible
    and reduce unnecessary rerouting after a network change.
\end{itemize}

The required response is strict JSON:

\begin{verbatim}
{"strategy":"protect_capacity",
 "reason":"Several affected flows share high-pressure links."}
\end{verbatim}

The LLM does not generate a route. Its response selects which deterministic ordering
and candidate-ranking rules the allocator will apply. Inference runs in a background
worker, so the simulation remains non-blocking. While a request is pending, an existing
valid plan can continue to operate. If no valid plan exists yet, the agent only retains
safe routine routes and leaves unresolved complex flows for a later cycle.

\subsection{Candidate Generation and Joint Allocation}

Candidate paths are generated deterministically with NetworkX. Several edge-cost views
are used---latency, residual capacity, utilization, and a balanced combination---and
duplicate paths are removed. Candidates that violate the exact end-to-end latency SLA
are discarded before allocation.

The \texttt{route\_with\_strategy} node applies the active LLM objective to all active
flows. The objective changes both the order in which flows are processed and the rank
of their candidate paths. The allocator then performs the following steps:

\begin{enumerate}
    \item generate feasible candidate paths for each active flow;
    \item preserve the current path as an option when it remains valid;
    \item order flows and rank candidates according to the selected objective;
    \item test each candidate against a shared capacity ledger;
    \item select the first feasible route and reserve its capacity;
    \item continue with the next flow.
\end{enumerate}

The allocator is therefore joint, sequential, greedy, and capacity-aware. It is
``joint'' because every decision shares one view of capacity, not because it solves an
exact global optimization problem. Consequently, its result can still depend on flow
ordering; the LLM objective influences that ordering and is therefore operationally
relevant.

\subsection{Validate}

The \texttt{validate} node is the final safety gate. Before a route reaches the
controller, this node checks path continuity, source and destination, link existence,
link availability, nominal capacity, the flow's exact maximum-latency SLA, and aggregate
capacity under the shared ledger. Invalid routes are removed.

This early filtering also affects how rejections appear in the results. A heuristic may
install a route and later count it as a latency violation, whereas the agent rejects the
same infeasible candidate before installation and records it as having no valid path.
Therefore, rejection categories must be interpreted together with total accepted
traffic rather than in isolation.

\subsection{Record}

The \texttt{record} node stores the situation, selected strategy, decision source, LLM
provider and model, inference time, token use, routed and unresolved flows, and invalid
actions. Each experiment exports \texttt{network\_metrics.csv},
\texttt{link\_history.csv}, \texttt{flow\_events.csv},
\texttt{model\_metrics.csv}, \texttt{agent\_decisions.csv}, and
\texttt{run\_metadata.json} in a model-specific directory. This prevents one experiment
from overwriting another and supplies the data used by the dashboard.


\section{Simulation Methodology}

\subsection{Models and baselines}

The final agent was evaluated with the following models:

\begin{itemize}
    \item OpenAI GPT-5.4 Nano;
    \item Anthropic Claude Haiku 4.5;
    \item Anthropic Claude Sonnet 4.5;
    \item Google Gemini 2.5 Flash;
    \item local Llama 3.2 1B running on CPU through Ollama.
\end{itemize}

DeepSeek and Groq configurations were explored during development, but were excluded
from the final comparison because their available endpoints did not produce sufficiently
consistent structured responses under the experiment configuration. The network
performance baselines were \texttt{heuristic\_no\_delay},
\texttt{heuristic\_low\_delay}, and \texttt{heuristic\_high\_delay}.

The recorded experiments used a random seed of 50 and 120 simulation steps. Most remote
model runs used Top-$K=10$; the recorded Sonnet and local Llama runs used Top-$K=5$.
This difference, together with asynchronous inference, must be considered when comparing
models. In particular, the Sonnet directory contains only the agentic run and therefore
does not provide a paired heuristic baseline from the same execution.

\subsection{Metrics}

Network performance was evaluated using:

\begin{itemize}
    \item packet delivery or traffic-acceptance ratio (PDR proxy), computed as accepted
    Mbps divided by offered Mbps;
    \item total delivered Mbps;
    \item number of accepted and rejected flow requests;
    \item latency-SLA violations and requests with no valid path;
    \item link utilization and capacity compliance.
\end{itemize}

Agent performance was evaluated using:

\begin{itemize}
    \item number of LLM requests and successful structured responses;
    \item input, output, reasoning, and total token usage;
    \item inference time;
    \item selected-strategy distribution and duration of each active plan;
    \item number of routed, unresolved, and invalid actions.
\end{itemize}

It is important to distinguish traffic acceptance from request acceptance. The former
weights flows by Mbps, whereas the latter gives every flow request the same weight.
Both are useful, but they answer different questions.


\section{Results and Discussion}

An interactive dashboard was implemented with Streamlit to compare all models and to
inspect each run individually. The dashboard combines network outcomes, link history,
flow events, model metrics, and agent decisions.

\subsection{Network performance}

Figure~\ref{fig:pdr-comparison} shows that the remote-model agentic configurations
outperformed the heuristic baselines in delivered traffic. The best result was obtained
with GPT-5.4 Nano, with a traffic-acceptance ratio of 38.56\%. Claude Sonnet and Claude
Haiku were almost identical at 38.50\%, while Gemini achieved 37.16\%. The differences
among the four remote models were therefore small: GPT exceeded Sonnet by only
0.06 percentage points and Haiku by approximately the same amount.

Within the paired GPT run, the best heuristic achieved 26.71\%, compared with 38.56\%
for the agentic configuration. This is an absolute improvement of 11.85 percentage
points, or a 44.4\% relative increase in accepted traffic. The paired improvements were
also positive for Haiku and Gemini. The local Llama run was the exception: it achieved
19.79\%, below its best paired heuristic result of 31.04\%.

\begin{figure}[H]
    \centering
    \includegraphics[width=0.6\linewidth]{Imagenes/pdr.png}
    \caption{Traffic-acceptance ratio over time for the agentic models and heuristic baselines.}
    \label{fig:pdr-comparison}
\end{figure}

The decrease near the end of the simulation is expected from the workload. New flows
continue to arrive and are not removed during the evaluated horizon, so offered demand
eventually exceeds the physical capacity of the topology. At that point, no routing
policy can accept every request. The relevant result is therefore not the existence of
the final decline, but that the agentic configurations sustain a higher delivered load
before and during saturation.

The request-level results support the same conclusion. Out of 12,784 requests, GPT
accepted 5,417 (42.37\%), Haiku accepted 5,404 (42.27\%), Sonnet accepted 5,408
(42.30\%), and Gemini accepted 5,210 (40.76\%). The local model accepted only 2,760
(21.59\%). These values differ from the PDR percentages because request acceptance is
not weighted by flow bandwidth.

Figure~\ref{fig:delivered-mbps} presents cumulative delivered traffic across the 120
steps. GPT delivered 41,005.60 Mbps in the aggregated per-step measurements. Sonnet
delivered 40,941.42 Mbps, Haiku 40,937.10 Mbps, and Gemini 39,514.98 Mbps. Once again,
the remote models produced similar results, whereas the local model reached only
21,045.21 Mbps.

\begin{figure}[H]
    \centering
    \includegraphics[width=0.6\linewidth]{Imagenes/entregados.png}
    \caption{Aggregated delivered Mbps over the simulation horizon.}
    \label{fig:delivered-mbps}
\end{figure}

Most agentic rejections were classified as \emph{no valid path}, while only two
latency-violation drops were recorded in each successful remote-model run. This does not
mean that latency was ignored. On the contrary, candidate paths that violated the SLA
were removed before installation. As a result, many requests that a heuristic might
attempt and later count as delay violations were rejected earlier by the agent. The
combination of a higher acceptance ratio and very few post-installation latency drops
indicates that the agent was not increasing throughput merely by violating the SLA.

\subsection{Contribution of the LLM}

The strategy logs confirm that the models did not always make the same decision. For
example, GPT predominantly used \texttt{protect\_capacity} and occasionally
\texttt{minimize\_delay}; Haiku used \texttt{protect\_capacity},
\texttt{minimize\_delay}, and \texttt{maximize\_acceptance}; Sonnet divided its active
cycles mainly between \texttt{maximize\_acceptance} and
\texttt{protect\_capacity}; and Gemini also used \texttt{stable\_recovery}. These
counts describe how long an active plan influenced the allocator, not only the number
of LLM calls.

Nevertheless, the small performance difference among remote models shows that the
deterministic part of the architecture contributes strongly to robustness. Capacity
checking, SLA filtering, candidate generation, and joint allocation constrain the
space in which the LLM can act. Several strategies can also converge to similar paths
when the same bottleneck dominates the network. The correct interpretation is therefore
that the LLM has a real but bounded influence: it changes flow ordering and candidate
ranking, while deterministic tools guarantee feasibility.

The current results alone cannot isolate exactly how much of the improvement is caused
by the LLM. A fixed-strategy ablation using the same allocator, seeds, and event sequence
would be required for a causal comparison. This is a more defensible conclusion than
attributing the complete agentic gain either to the model or to the allocator.

\subsection{Token usage and inference latency}

Figure~\ref{fig:token-usage} shows that Claude Haiku consumed the most tokens: 21,748
across 15 completed calls. GPT completed 17 calls using 11,302 tokens, Sonnet completed
9 calls using 11,565 tokens, and Gemini completed 5 calls using 7,597 tokens. Token
usage is not determined only by model size. It also depends on response style, reasoning
tokens, prompt size, call frequency, and how long an asynchronous request remains in
flight. Gemini, for example, used substantial reasoning tokens per response but made
fewer completed calls because each call took longer.

\begin{figure}[H]
    \centering
    \includegraphics[width=0.6\linewidth]{Imagenes/tokens.png}
    \caption{Total token usage by model.}
    \label{fig:token-usage}
\end{figure}

GPT was the fastest remote model, with a mean inference time of 1.38 s, followed by
Haiku at 1.69 s, Sonnet at 2.61 s, and Gemini at 5.09 s. The local 1B model recorded one
completion after 16.94 s. Its response echoed part of the prompt rather than returning
a usable strategy object, so no valid active plan was installed. It is therefore more
accurate to say that the local model produced one nominal completion but no effective
routing decision, rather than saying that it never returned anything.

\begin{figure}[H]
    \centering
    \includegraphics[width=0.6\linewidth]{Imagenes/inference.png}
    \caption{Mean LLM inference time by model.}
    \label{fig:inference-time}
\end{figure}

The asynchronous design prevented these inference times from directly blocking the
simulation. However, slower inference still has an indirect cost: a plan arrives later,
can influence fewer cycles, and has a greater probability of being based on an older
snapshot. Event-driven calls and active-plan reuse reduce this problem but do not remove
it completely.

\subsection{Threats to validity}

The experiments demonstrate promising behavior, but several limitations should be
considered. First, a single random seed and one run per configuration are not sufficient
to establish statistical significance. Second, Top-$K$ was not identical in every
recorded run, and Sonnet did not have a paired heuristic run. Third, asynchronous model
latency changes the number and timing of completed decisions, so model quality and model
speed are partially coupled. Finally, the heuristic results vary across some saved run
directories; comparisons should therefore use paired results from the same run whenever
possible. Repeated experiments with identical configurations, multiple seeds, confidence
intervals, and a fixed-strategy ablation would strengthen the evaluation.


\section*{Conclusion and Future Work}

The project produced a working LangGraph routing agent that combines event-driven LLM
reasoning with deterministic network control. The final architecture addresses the
main failures observed in earlier versions: it avoids per-flow LLM scaling, does not
call the model every cycle, prevents independent capacity allocation, validates hard
constraints before installation, and records sufficient telemetry for model comparison.

The remote-model configurations delivered substantially more traffic than their paired
heuristic baselines while producing almost no post-installation latency violations. GPT,
Haiku, and Sonnet achieved nearly identical network performance, with Gemini slightly
behind. The local CPU model was too slow and did not produce a usable structured plan.

The results also reveal the intended trade-off of the architecture. The LLM influences
the routing objective, but its freedom is constrained by deterministic candidate
generation, shared-capacity allocation, and SLA validation. This makes the system safer
and more stable, although it also makes the model's isolated contribution difficult to
measure.

Future work should first perform controlled ablations in which the same allocator uses
a fixed objective, a random objective, and the LLM-selected objective. Experiments should
then be repeated over multiple seeds. Other useful extensions include an exact or
look-ahead allocator to compare against the current greedy method, adaptive Top-$K$
selection, explicit detection and rejection of stale LLM responses, and carefully
designed short-term reflection based on measurable action--outcome pairs. Memory should
only be reintroduced after the system can attribute an outcome to a specific decision;
otherwise, it risks preserving noise rather than learning from experience.
